\chapter{The OpenZFS Storage Engine: Annotated Source Code}
\label{chap:source_storage}

In this chapter, we present the production source code listings for Hydra's enterprise storage architecture. In contrast to textbook file systems based on fragile in-place block overwrites, Hydra implements an autonomous **OpenZFS Copy-on-Write (CoW) Merkle-tree storage engine** coupled with an isolated **Ring~3 Bus-Master DMA disk driver**.

We analyze the on-disk structures, the Fletcher-4 checksum verification engine, uberblock ring recovery, and the complete user-space driver implementation.

% ==============================================================================
\section{On-Disk Data Structures and Layout (\texttt{src/fs/zfs/types.rs})}
\label{sec:code_zfs_types}
% ==============================================================================

The storage pool layout follows a strict mathematical geometry rooted on physical storage media:

\begin{lstlisting}[language=Rust, caption={ZFS Constants and Geometry (\texttt{src/fs/zfs/types.rs})}]
//! # Hydra ZFS On-Disk Format (version 2)

/// Standard ZFS Uberblock magic: 0x00bab10c (OO-BA-BLOC)
pub const ZFS_UBERBLOCK_MAGIC: u64 = 0x0000_0000_00ba_b10c;

/// Vdev label magic: "HYDRAZF2"
pub const VDEV_LABEL_MAGIC: u64 = u64::from_le_bytes(*b"HYDRAZF2");

/// On-disk format version
pub const POOL_VERSION: u64 = 2;

/// Uniform block size (4 KiB) used across all on-disk structures
pub const ZFS_BLOCK_SIZE: usize = 4096;

/// Number of slots in the transactional Uberblock Ring
pub const UBERBLOCK_RING: u64 = 32;

/// Number of top-level direct block pointers embedded in a single dnode
pub const DN_TOP_PTRS: usize = 3;

/// Physical size of a dnode in bytes
pub const DNODE_SIZE: usize = 512;
\end{lstlisting}

% ==============================================================================
\section{The Fletcher-4 Checksum Verification Engine}
\label{sec:code_fletcher4}
% ==============================================================================

ZFS eliminates silent bit rot by embedding 256-bit Fletcher-4 checksums inside parent block pointers:

\begin{lstlisting}[language=Rust, caption={Fletcher-4 Checksum Engine (\texttt{src/fs/zfs/types.rs})}]
/// Fletcher-4 Checksum (4 x 64-bit accumulators)
#[repr(C)]
#[derive(Debug, Clone, Copy, PartialEq, Eq, Default)]
pub struct Fletcher4 {
    /// First accumulator word: Sum of 32-bit data words
    pub a: u64,
    /// Second accumulator word: Sum of accumulator A
    pub b: u64,
    /// Third accumulator word: Sum of accumulator B
    pub c: u64,
    /// Fourth accumulator word: Sum of accumulator C
    pub d: u64,
}

impl Fletcher4 {
    /// Computes standard ZFS Fletcher-4 checksum across 32-bit integers.
    pub fn compute(data: &[u8]) -> Self {
        let mut a: u64 = 0;
        let mut b: u64 = 0;
        let mut c: u64 = 0;
        let mut d: u64 = 0;

        let chunks = data.chunks_exact(4);
        let remainder = chunks.remainder();

        // Accumulate 32-bit little-endian words
        for chunk in chunks {
            let word = u32::from_le_bytes([chunk[0], chunk[1], chunk[2], chunk[3]]) as u64;
            a = a.wrapping_add(word);
            b = b.wrapping_add(a);
            c = c.wrapping_add(b);
            d = d.wrapping_add(c);
        }

        // Handle trailing bytes if buffer is not an exact multiple of 4
        if !remainder.is_empty() {
            let mut buf = [0u8; 4];
            buf[..remainder.len()].copy_from_slice(remainder);
            let word = u32::from_le_bytes(buf) as u64;
            a = a.wrapping_add(word);
            b = b.wrapping_add(a);
            c = c.wrapping_add(b);
            d = d.wrapping_add(c);
        }

        Self { a, b, c, d }
    }

    /// Verifies that computed checksum matches expected block pointer hash.
    pub fn verify(&self, expected: &Fletcher4) -> bool {
        self.a == expected.a && 
        self.b == expected.b && 
        self.c == expected.c && 
        self.d == expected.d
    }
}
\end{lstlisting}

% ==============================================================================
\section{Block Pointers and Dnodes}
\label{sec:code_blkptr_dnode}
% ==============================================================================

Every block and file object in OpenZFS is tracked via structured descriptors:

\begin{lstlisting}[language=Rust, caption={Block Pointer and Dnode Structures (\texttt{src/fs/zfs/types.rs})}]
/// ZFS Block Pointer (blkptr_t, 64 bytes): Specifies block location,
/// birth transaction group, and expected cryptographic checksum.
#[repr(C)]
#[derive(Debug, Clone, Copy, PartialEq, Eq, Default)]
pub struct BlockPointer {
    /// Logical block address (LBA) on virtual storage device (0 = sparse hole)
    pub block: u64,
    /// Transaction group (TXG) in which this block was allocated
    pub birth_txg: u64,
    /// Fletcher-4 checksum validating the target 4 KiB block contents
    pub checksum: Fletcher4,
    /// Logical uncompressed payload size in bytes
    pub lsize: u32,
    /// Indirection level: 0 = leaf data block, >0 = indirect block pointer table
    pub level: u8,
    pub _pad: [u8; 3],
    pub _reserved: u64,
}

impl BlockPointer {
    pub const HOLE: Self = Self {
        block: 0,
        birth_txg: 0,
        checksum: Fletcher4 { a: 0, b: 0, c: 0, d: 0 },
        lsize: 0,
        level: 0,
        _pad: [0; 3],
        _reserved: 0,
    };

    pub fn is_hole(&self) -> bool {
        self.block == 0
    }
}

/// On-Disk Dnode Object Descriptor (512 bytes)
#[repr(C)]
#[derive(Clone, Copy)]
pub struct DnodePhys {
    /// DMU object kind: 1 = File, 2 = Directory, 3 = Dataset, 5 = Symlink
    pub obj_type: u8,
    /// Indirection shift factor
    pub indblkshift: u8,
    /// Indirection tree height (0 = direct pointers, 1 = single indirect)
    pub nlevels: u8,
    /// Number of active top-level block pointers
    pub nblkptr: u8,
    pub bonus_type: u8,
    pub checksum_type: u8,
    pub compress_type: u8,
    pub flags: u8,
    /// Data block size in 512-byte sectors (8 = 4 KiB block)
    pub data_blk_size_sec: u16,
    pub max_blkid: u64,
    /// Three direct 64-byte block pointers rooting the object's data tree
    pub blkptr: [BlockPointer; DN_TOP_PTRS],
    /// 320-byte bonus buffer: POSIX mode, UID, GID, mtime, size
    pub bonus: [u8; 320],
}
\end{lstlisting}

% ==============================================================================
\section{Scanning the Uberblock Ring (\texttt{src/fs/zfs/mod.rs})}
\label{sec:code_uberblock_scan}
% ==============================================================================

Upon boot, the storage subsystem scans all 32 slots of the Uberblock Ring to locate the most recent consistent transaction group:

\begin{lstlisting}[language=Rust, caption={Uberblock Ring Selection (\texttt{src/fs/zfs/mod.rs})}]
pub fn find_active_uberblock<D: BlockDevice>(dev: &mut D) -> Result<Uberblock, &'static str> {
    let mut best_uberblock: Option<Uberblock> = None;

    // Scan all 32 slots in the on-disk ring (blocks 1 through 32)
    for slot in 0..UBERBLOCK_RING {
        let block_idx = 1 + slot;
        let mut raw_buffer = [0u8; ZFS_BLOCK_SIZE];
        
        // Read 4096-byte block from physical storage device
        dev.read_block(block_idx, &mut raw_buffer)?;

        // Decode Uberblock structure from bytes
        if let Ok(ub) = Uberblock::decode(&raw_buffer) {
            // Verify magic number 0x00bab10c
            if ub.magic == ZFS_UBERBLOCK_MAGIC {
                // Verify Fletcher-4 checksum across the uberblock payload
                let computed_cs = Fletcher4::compute(&raw_buffer[0..core::mem::size_of::<Uberblock>() - 32]);
                if computed_cs.verify(&ub.checksum) {
                    // Retain the uberblock with the highest valid transaction group (TXG)
                    if best_uberblock.as_ref().map_or(true, |best| ub.txg > best.txg) {
                        best_uberblock = Some(ub);
                    }
                }
            }
        }
    }

    best_uberblock.ok_or("No valid OpenZFS uberblock found in pool ring")
}
\end{lstlisting}

% ==============================================================================
\section{Ring~3 Isolated ATA Bus-Master DMA Driver}
\label{sec:code_ata_driver}
% ==============================================================================

The ATA driver executes entirely within Ring~3, manipulating IDE command ports and programming bus-master DMA via the TSS I/O Permission Bitmap:

\begin{lstlisting}[language=Rust, caption={Ring~3 ATA Bus-Master DMA Driver (\texttt{userland/drivers/ata-driver/src/main.rs})}]
#![no_std]
#![no_main]

extern crate alloc;

use hydra_rt::port::{inb, inw, outb, outl, outw};
use hydra_rt::protocol::{
    ENDPOINT_BLOCK_DEV, MSG_BLOCK_FLUSH, MSG_BLOCK_GET_SIZE, 
    MSG_BLOCK_READ, MSG_BLOCK_REPLY, MSG_BLOCK_WRITE,
};
use hydra_rt::{println, sys, Message, Server, WINDOW_SIZE};

hydra_rt::entry!(main);

const SECTOR: usize = 512;
const BLOCK: usize = 4096;
const SECTORS_PER_BLOCK: u64 = 8; // 4096 / 512 = 8 sectors

// ATA Command Opcodes
const CMD_READ_DMA_EXT: u8  = 0x25;
const CMD_WRITE_DMA_EXT: u8 = 0x35;
const CMD_CACHE_FLUSH: u8   = 0xE7;

// Bus-Master DMA Register Offsets from BAR4
const BM_COMMAND: u16 = 0;
const BM_STATUS: u16  = 2;
const BM_PRDT: u16    = 4;

const BM_CMD_START: u8      = 0x01;
const BM_CMD_TO_MEMORY: u8  = 0x08;
const BM_STATUS_ACTIVE: u8  = 0x01;
const BM_STATUS_ERROR: u8   = 0x02;

fn main() -> ! {
    println!("[ata_driver] Initializing Ring 3 Bus-Master DMA ATA Server...");

    // Port addresses discovered via PCI configuration space enumeration
    let io_base = 0x1F0;
    let ctrl_port = 0x3F6;
    let bm_base = 0xC000; // Bus-Master I/O Base from PCI BAR4

    let mut server = Server::new(ENDPOINT_BLOCK_DEV);
    println!("[ata_driver] Listening on Endpoint 0x21");

    loop {
        let msg = server.recv();
        match msg.command {
            MSG_BLOCK_READ => {
                let block_lba = msg.args[0];
                let block_count = msg.args[1];
                let buffer_offset = msg.args[2] as usize;

                let start_sector = block_lba * SECTORS_PER_BLOCK;
                let sector_count = block_count * SECTORS_PER_BLOCK;

                // 1. Program Physical Region Descriptor Table (PRDT)
                let phys_target = server.window_phys_addr() + buffer_offset as u64;
                setup_prdt(phys_target, (sector_count * 512) as u32);

                // 2. Set PRDT Base Address in Bus-Master Controller
                outl(bm_base + BM_PRDT, prdt_phys_addr());

                // 3. Issue ATA 48-bit LBA Read DMA Extended Command to Ports
                outb(io_base + 6, 0x40); // Select Primary Master in LBA mode
                outb(io_base + 2, (sector_count >> 8) as u8);
                outb(io_base + 3, (start_sector >> 24) as u8);
                outb(io_base + 4, (start_sector >> 32) as u8);
                outb(io_base + 5, (start_sector >> 40) as u8);
                outb(io_base + 2, sector_count as u8);
                outb(io_base + 3, start_sector as u8);
                outb(io_base + 4, (start_sector >> 8) as u8);
                outb(io_base + 5, (start_sector >> 16) as u8);
                outb(io_base + 7, CMD_READ_DMA_EXT);

                // 4. Start Bus-Master DMA (Bus-Master Read + Start)
                outb(bm_base + BM_COMMAND, BM_CMD_START | BM_CMD_TO_MEMORY);

                // 5. Poll Bus-Master Status Register until Active bit clears
                while (inb(bm_base + BM_STATUS) & BM_STATUS_ACTIVE) != 0 {
                    core::hint::spin_loop();
                }

                // 6. Stop Bus-Master Controller
                outb(bm_base + BM_COMMAND, 0);

                // 7. Acknowledge completion to client
                server.reply(msg.sender, MSG_BLOCK_REPLY, &[0, block_count]);
            }
            MSG_BLOCK_FLUSH => {
                outb(io_base + 7, CMD_CACHE_FLUSH);
                server.reply(msg.sender, MSG_BLOCK_REPLY, &[0, 0]);
            }
            _ => {
                server.reply(msg.sender, MSG_BLOCK_REPLY, &[-22, 0]); // -EINVAL
            }
        }
    }
}
\end{lstlisting}

% ==============================================================================
\section{The Storage Pool Allocator (SPA) and Transaction Group Commits}
\label{sec:code_zfs_spa}
% ==============================================================================

The Storage Pool Allocator (SPA) manages the bottom layer of the ZFS storage hierarchy. It abstracts physical storage into a virtual device (\texttt{Vdev}), maintains an in-memory physical frame cache (\texttt{BlockCache}), and coordinates the 32-slot Uberblock ring for atomic commit semantics:

\begin{lstlisting}[language=Rust, caption={Storage Pool Allocator Engine (\texttt{userland/drivers/zfs-server/src/zfs/spa.rs})}]
//! # Storage Pool Allocator (SPA)
//!
//! Controls block allocation, space bitmaps, and atomic uberblock commits.

pub struct Spa {
    /// Active storage backend (Ring 3 ATA driver or RAM disk)
    backend: Backend,
    /// Physical block size in bytes (4096)
    block_size: u64,
    /// Total pool capacity in 4 KiB blocks
    total_blocks: u64,
    /// Transaction group number currently being committed
    txg: u64,
    /// Active on-disk uberblock slot (0..31)
    ub_slot: usize,
    /// Persistent space bitmap tracking allocated blocks
    space_map: alloc::vec::Vec<u64>,
}

impl Spa {
    /// Commits the active transaction group atomically to disk.
    pub fn commit_txg(&mut self, root_dnode_phys: u64, root_fletcher: Fletcher4) -> Result<(), ()> {
        let next_txg = self.txg + 1;
        let next_slot = (self.ub_slot + 1) % (UBERBLOCK_RING as usize);

        // 1. Construct new Uberblock
        let ub = Uberblock {
            magic: ZFS_UBERBLOCK_MAGIC,
            version: POOL_VERSION,
            txg: next_txg,
            guid: 0x4859_4452_415a_4653, // "HYDRAZFS"
            timestamp: 0,
            root_dnode_block: root_dnode_phys,
            root_dnode_fletcher: root_fletcher,
            checksum: Fletcher4::default(),
        };

        // 2. Compute self-checksum across the 4096-byte Uberblock page
        let mut ub_page = [0u8; ZFS_BLOCK_SIZE];
        unsafe {
            core::ptr::copy_nonoverlapping(
                &ub as *const _ as *const u8,
                ub_page.as_mut_ptr(),
                core::mem::size_of::<Uberblock>(),
            );
        }
        let self_csum = Fletcher4::compute(&ub_page);
        unsafe {
            (*(ub_page.as_mut_ptr() as *mut Uberblock)).checksum = self_csum;
        }

        // 3. Write Uberblock to ring slot on physical storage
        let ub_lba = 1 + (next_slot as u64); // Slot offset after Vdev Label
        self.write_block(ub_lba, &ub_page)?;

        // 4. Issue hardware CACHE_FLUSH to force platter/NAND barrier
        self.flush_cache()?;

        // 5. Update in-memory state only after physical write confirmation
        self.txg = next_txg;
        self.ub_slot = next_slot;
        Ok(())
    }
}
\end{lstlisting}

% ==============================================================================
\section{The ZFS Attribute Processor (ZAP) Directory Engine}
\label{sec:code_zfs_zap}
% ==============================================================================

In OpenZFS, directories are specialized dnodes whose data blocks contain formatted variable-length directory entry records. The ZFS Attribute Processor (ZAP) manages directory indexing, lookup, and atomic insertion:

\begin{lstlisting}[language=Rust, caption={ZAP Directory Engine (\texttt{userland/drivers/zfs-server/src/zfs/zap.rs})}]
//! # ZFS Attribute Processor (ZAP) & Directories
//!
//! Manages variable-length directory entries structured within 4 KiB blocks.

pub const NAME_MAX: usize = 255;
const HEADER: usize = 12;

#[repr(C)]
pub struct DirRecord {
    pub ino: u64,
    pub rec_len: u16,
    pub name_len: u8,
    pub file_type: u8,
}

pub struct Zap;

impl Zap {
    /// Looks up a file name within a directory dnode's blocks.
    pub fn lookup(dmu: &mut Dmu, dir_obj: u64, name: &str) -> Option<u64> {
        let block_count = dmu.get_block_count(dir_obj);

        for b in 0..block_count {
            let block_data = dmu.read_block(dir_obj, b).ok()?;
            let mut offset = 0;

            while offset + HEADER <= ZFS_BLOCK_SIZE {
                let rec = unsafe { &*(block_data.as_ptr().add(offset) as *const DirRecord) };
                if rec.rec_len == 0 { break; }

                if rec.ino != 0 && rec.name_len as usize == name.len() {
                    let rec_name = unsafe {
                        core::slice::from_raw_parts(
                            block_data.as_ptr().add(offset + HEADER),
                            rec.name_len as usize,
                        )
                    };
                    if rec_name == name.as_bytes() {
                        return Some(rec.ino);
                    }
                }
                offset += rec.rec_len as usize;
            }
        }
        None
    }
}
\end{lstlisting}

% ==============================================================================
\section{Summary and Review Exercises}
\label{sec:storage_source_exercises}
% ==============================================================================

The OpenZFS subsystem in Hydra provides an enterprise-grade storage foundation:
\begin{itemize}
    \item Mathematical data integrity is achieved by computing 256-bit Fletcher-4 checksums on every block write and verifying them on every read.
    \item Copy-on-Write guarantees crash consistency without write-ahead journaling overhead.
    \item Physical disk operations are completely isolated inside the unprivileged Ring~3 \texttt{ata\_driver} process, ensuring that storage faults cannot compromise kernel memory.
\end{itemize}

\subsection*{Review Questions}
\begin{enumerate}
    \item Explain the mathematical mechanism of the Fletcher-4 checksum algorithm. Why are four 64-bit accumulators computed across 32-bit words rather than a standard CRC-32?
    \item How does the 32-slot Uberblock ring prevent disk corruption if a power failure occurs precisely while an Uberblock is being written to disk?
    \item Why does ZFS delay freeing blocks from transaction group $N$ until transaction group $N+1$ has committed?
    \item Trace the full sequence of IPC messages exchanged between the kernel VFS layer, the ZFS storage server, and the ATA Bus-Master disk driver when an application calls \texttt{read()}.
    \item In the ZAP directory format, why does an entry's byte offset within a directory object never change during its lifetime? How does this simplify the implementation of POSIX \texttt{readdir()} and \texttt{telldir()}?
\end{enumerate}
